\section{Approximation simplifiée du spectre angulaire $C_\ell$}

\subsection{Expression exacte}

Le spectre de puissance angulaire des fluctuations projetées s'écrit, dans une approximation de type Limber :

\begin{equation}
C_\ell = \int dz ; \frac{H(z)}{r^2(z)} ; n^2(z); P!\left(k=\frac{\ell}{r(z)}, z\right),
\end{equation}

où :
\begin{itemize}
\item $n(z)$ est la distribution normalisée en redshift,
\item $r(z)$ est la distance comobile,
\item $H(z)$ est le taux d'expansion,
\item $P(k,z)$ est le spectre de puissance 3D de la matière.
\end{itemize}

Cette expression est celle calculée numériquement par des codes comme \texttt{CCL}.

---

\subsection{Objectif de la simplification}

Nous cherchons à comprendre comment le dipôle de clustering (i.e. $C_1$) dépend de la distribution en redshift $n(z)$, en particulier sous un cut $z > z_{\min}$.

L'objectif n'est pas de reproduire précisément $C_\ell$, mais d'isoler sa dépendance principale en $n(z)$.

---

\subsection{Approximation 1 : séparation du spectre de puissance}

On suppose que le spectre de puissance peut être séparé en une dépendance en redshift et en échelle :

\begin{equation}
P(k,z) \simeq D^2(z), P(k),
\end{equation}

où $D(z)$ est le facteur de croissance linéaire.

Si $D(z)$ varie lentement sur le support de $n(z)$, on peut l'absorber dans une constante de normalisation :

\begin{equation}
P(k,z) \approx P(k) \times \text{const}.
\end{equation}

---

\subsection{Approximation 2 : échelle effective}

L'intégrale contient une dépendance en $k = \ell / r(z)$.
On approxime cette dépendance en introduisant une distance effective $r_{\rm eff}$ telle que :

\begin{equation}
k \approx \frac{\ell}{r_{\rm eff}}.
\end{equation}

Cette approximation revient à supposer que la projection est dominée par une échelle caractéristique.

On obtient alors :

\begin{equation}
C_\ell \propto P!\left(\frac{\ell}{r_{\rm eff}}\right)
\int dz ; \frac{n^2(z)}{r^2(z)}.
\end{equation}

---

\subsection{Approximation 3 : forme du spectre de puissance}

À grande échelle (petits $k$), le spectre de puissance de la matière suit approximativement une loi de puissance :

\begin{equation}
P(k) \propto k^{n_{\rm eff}},
\end{equation}

avec un indice effectif typique :

\begin{equation}
n_{\rm eff} \approx -1.2.
\end{equation}

Ce comportement provient de la combinaison :
\begin{itemize}
\item du spectre primordial ($n_s \simeq 1$),
\item de la fonction de transfert (suppression aux grandes échelles),
\item de la croissance des structures.
\end{itemize}

---

\subsection{Passage à la dépendance en $\ell$}

En remplaçant $k = \ell / r_{\rm eff}$, on obtient :

\begin{equation}
P!\left(\frac{\ell}{r_{\rm eff}}\right)
\propto \left(\frac{\ell}{r_{\rm eff}}\right)^{n_{\rm eff}}
= r_{\rm eff}^{-n_{\rm eff}} , \ell^{n_{\rm eff}}.
\end{equation}

Le facteur $r_{\rm eff}^{-n_{\rm eff}}$ étant constant, il peut être absorbé dans la normalisation globale.

Ainsi :

\begin{equation}
C_\ell \propto \ell^{n_{\rm eff}} ;
\int dz ; \frac{n^2(z)}{r^2(z)}.
\end{equation}

En utilisant $n_{\rm eff} \approx -1.2$, on obtient :

\begin{equation}
C_\ell \propto \ell^{-1.2} ;
\int dz ; \frac{n^2(z)}{r^2(z)}.
\end{equation}

---

\subsection{Interprétation physique}

L'expression obtenue montre que :

\begin{itemize}
\item la \textbf{forme en $\ell$} est fixée par le spectre de puissance 3D,
\item l'\textbf{amplitude} est entièrement contrôlée par :
\begin{equation}
\int dz ; \frac{n^2(z)}{r^2(z)}.
\end{equation}
\end{itemize}

Cette intégrale a une interprétation claire :
\begin{itemize}
\item le terme $n^2(z)$ provient du fait que l'on corrèle des paires d'objets,
\item le facteur $1/r^2(z)$ traduit la dilution géométrique des structures projetées.
\end{itemize}

---

\subsection{Conséquence pour le dipôle}

Le dipôle de clustering est relié à $C_1$ par :

\begin{equation}
d_{\rm RMS} = \sqrt{\frac{3C_1}{4\pi}}.
\end{equation}

Ainsi, toute variation de $n(z)$ (par exemple un cut $z > z_{\min}$) agit directement sur :

\begin{equation}
C_1 \propto \int dz ; \frac{n^2(z)}{r^2(z)}.
\end{equation}

En particulier, supprimer les faibles redshifts (petits $r$) réduit fortement cette intégrale, ce qui diminue rapidement le dipôle induit par le clustering.

---

\subsection{Résumé}

Le modèle simplifié repose donc sur la factorisation :

\begin{equation}
C_\ell \approx \left[\int dz ; \frac{n^2(z)}{r^2(z)}\right] \times \ell^{-1.2},
\end{equation}

qui capture correctement la dépendance du dipôle de clustering en fonction de la distribution en redshift, tout en restant beaucoup plus simple que le calcul complet.
